\documentclass{article}
\usepackage{amsmath,amssymb,amsthm}
\usepackage{algorithm,algorithmic}
\usepackage{hyperref}
\usepackage{bm}
\usepackage{enumitem}
\usepackage{geometry}
\geometry{margin=1in}
\newtheorem{theorem}{Theorem}
\newtheorem{lemma}{Lemma}
\newtheorem{assumption}{Assumption}
\newcommand{\EE}{\mathbb{E}}
\newcommand{\RR}{\mathbb{R}}
\newcommand{\norm}[1]{\left\|#1\right\|}
\newcommand{\inner}[2]{\langle #1,#2\rangle}
\begin{document}

\section*{Algorithm Name}
\textbf{Stochastic Gradient Langevin Dynamics with Decaying Noise (SGLD‑D)}  

The method combines a deterministic gradient step with an additive Gaussian perturbation whose variance is scheduled to decay over the iterations.

\section*{Mathematical Formulation}
Let $f:\RR^{d}\rightarrow\RR$ be the objective function.  
Given a stepsize (learning rate) $\alpha>0$ and a non‑increasing noise schedule $\{\tau_k\}_{k\ge 0}$, the iterates $\{x_k\}_{k\ge 0}$ are generated by  

\[
\boxed{
x_{k+1}=x_k-\alpha \nabla f(x_k)+\sqrt{2\alpha \tau_k}\,\xi_k,\qquad k=0,1,\dots
}
\tag{1}
\]

where $\xi_k\sim \mathcal N(0,I_d)$ are i.i.d. standard Gaussian vectors, independent of the history $\mathcal F_k:=\sigma(x_0,\xi_0,\dots,\xi_{k-1})$.  
The noise variance decays, for instance  

\[
\tau_k=\frac{\tau_0}{k+1},\qquad \tau_0>0. \tag{2}
\]

When $\tau_k\equiv 0$, (1) reduces to the classical Gradient Descent (GD).  

\section*{Pseudocode}
\begin{algorithm}[H]
\caption{SGLD‑D}
\begin{algorithmic}[1]
\REQUIRE Initial point $x_0\in\RR^d$, stepsize $\alpha>0$, initial temperature $\tau_0>0$
\FOR{$k=0,1,\dots,T-1$}
    \STATE Compute deterministic gradient $g_k:=\nabla f(x_k)$
    \STATE Sample $\xi_k\sim\mathcal N(0,I_d)$
    \STATE Set $\tau_k\gets \tau_0/(k+1)$
    \STATE Update 
    \[
    x_{k+1}\gets x_k-\alpha g_k+\sqrt{2\alpha \tau_k}\,\xi_k
    \]
\ENDFOR
\RETURN $x_T$
\end{algorithmic}
\end{algorithm}

\section*{Dry Run Example}
Consider the one‑dimensional quadratic  

\[
f(w)=(w-3)^2 .
\]

\begin{enumerate}
\item \textbf{Initialization:} $w_0=0$, stepsize $\alpha=0.1$, $\tau_0=1$.
\item \textbf{Iteration $k=0$:}
    \begin{itemize}
        \item Gradient $g_0=\nabla f(w_0)=2(w_0-3)=-6$.
        \item $\tau_0=\frac{1}{0+1}=1$, draw a sample $\xi_0\sim\mathcal N(0,1)$; suppose $\xi_0=0.5$.
        \item Update 
        \[
        w_1=0-0.1(-6)+\sqrt{2\cdot0.1\cdot1}\,0.5
        =0+0.6+ \sqrt{0.2}\,0.5\approx0.6+0.4472\cdot0.5\approx0.8236 .
        \]
        \item Objective $f(w_1)=(0.8236-3)^2\approx4.71$.
    \end{itemize}
\item \textbf{Iteration $k=1$:}
    \begin{itemize}
        \item $g_1=2(w_1-3)=2(0.8236-3)=-4.3528$.
        \item $\tau_1=\frac{1}{2}=0.5$, draw $\xi_1=-0.3$.
        \item Update 
        \[
        w_2=w_1-0.1g_1+\sqrt{2\cdot0.1\cdot0.5}\,(-0.3)
        =0.8236-0.1(-4.3528)-\sqrt{0.1}\,0.3
        \approx0.8236+0.4353-0.3162\cdot0.3\approx1.1315 .
        \]
        \item $f(w_2)=(1.1315-3)^2\approx3.48$.
    \end{itemize}
\item \textbf{Iteration $k=2$:}
    \begin{itemize}
        \item $g_2=2(w_2-3)=2(1.1315-3)=-3.7370$.
        \item $\tau_2=\frac{1}{3}\approx0.333$, draw $\xi_2=0.1$.
        \item Update 
        \[
        w_3=w_2-0.1g_2+\sqrt{2\cdot0.1\cdot0.333}\,0.1
        =1.1315-0.1(-3.7370)+\sqrt{0.0666}\,0.1
        \approx1.1315+0.3737+0.2582\cdot0.1\approx1.5200 .
        \]
        \item $f(w_3)=(1.5200-3)^2\approx2.18$.
    \end{itemize}
\end{enumerate}
The objective value decreases despite the injected noise, illustrating the balance between descent and exploration.

\section*{Assumptions}
\begin{assumption}[Convexity and Smoothness]\label{ass:convex}
The function $f$ is convex and $L$‑smooth, i.e. for all $x,y\in\RR^d$,
\[
f(y)\le f(x)+\inner{\nabla f(x)}{y-x}+\frac{L}{2}\norm{y-x}^2 .
\tag{3}
\]
\end{assumption}
\begin{assumption}[Existence of a Minimizer]\label{ass:opt}
There exists $x^\star\in\arg\min f$, and we denote $f^\star:=f(x^\star)$.
\end{assumption}
\begin{assumption}[Bounded Domain]\label{ass:bound}
All iterates remain in a bounded set $\mathcal X$ with diameter $D$, i.e. $\norm{x_k-x^\star}\le D$ for all $k$.
\end{assumption}
\begin{assumption}[Noise Schedule]\label{ass:noise}
The temperature sequence $\{\tau_k\}$ is non‑negative, non‑increasing and satisfies  
\[
\sum_{k=0}^{\infty}\tau_k <\infty ,\qquad 
\tau_k = \frac{\tau_0}{k+1}\;\; \text{for some }\tau_0>0 .
\tag{4}
\]
\end{assumption}

\section*{Main Convergence Theorem}
\begin{theorem}[Expected Suboptimality of SGLD‑D]\label{thm:main}
Under Assumptions \ref{ass:convex}–\ref{ass:noise} and a constant stepsize $\alpha\in\bigl(0,\tfrac{1}{L}\bigr]$, the iterates generated by \eqref{1} satisfy for any $K\ge 1$
\[
\EE\!\bigl[f(x_K)\bigr]-f^\star
\;\le\;
\frac{\norm{x_0-x^\star}^2}{2\alpha K}
\;+\;
\alpha L\,\tau_0\,\frac{\log (K+1)}{K}.
\tag{5}
\]
Consequently, the algorithm attains the optimal $O(1/K)$ rate up to a logarithmic factor induced by the decaying noise.
\end{theorem}

\section*{Theoretical Properties}
\begin{enumerate}[label=\arabic*.]
\item \textbf{Stability.} The recursion \eqref{1} is a contraction in expectation when $\alpha\le 1/L$, because the smoothness constant $L$ controls the growth of the gradient term.
\item \textbf{Hyperparameter Sensitivity.} The bound \eqref{5} shows a linear dependence on $\alpha$ in the bias term $\alpha L\tau_0\log(K+1)/K$ and an inverse dependence in the variance term $\norm{x_0-x^\star}^2/(2\alpha K)$. Hence the optimal constant stepsize balances these two contributions, yielding $\alpha^\star =\Theta\!\bigl(\sqrt{\frac{\norm{x_0-x^\star}^2}{L\tau_0\log(K+1)}}\bigr)$.
\item \textbf{Comparison to Baselines.}
    \begin{itemize}
        \item \emph{Gradient Descent (GD)} corresponds to $\tau_k\equiv0$; the bound reduces to $\frac{\norm{x_0-x^\star}^2}{2\alpha K}$, the classic $O(1/K)$ rate.
        \item \emph{Standard SGLD} with a fixed temperature $\tau_k\equiv\tau>0$ yields a steady-state bias $\alpha L\tau$, which does not vanish. The decaying schedule \eqref{4} eliminates this asymptotic bias, at the price of a $\log K/K$ factor.
    \end{itemize}
\end{enumerate}

\section*{Discussion}
The theorem demonstrates that, despite the stochastic perturbations, the iterates converge in expectation to the optimum at the same order as deterministic GD. The decreasing temperature plays the same role as a diminishing variance in classical stochastic approximation, guaranteeing that the cumulative noise remains finite (cf. \eqref{4}). The extra logarithmic term is unavoidable under the simple schedule \eqref{2}; more aggressive decay (e.g., $\tau_k\propto 1/(k+1)^{1+\epsilon}$) would remove the $\log$ factor at the expense of slower initial exploration.

\section*{Preliminaries (Definitions and Lemmas)}
\begin{lemma}[Smoothness Inequality]\label{lem:smooth}
If $f$ is $L$‑smooth then for any $x,y\in\RR^d$,
\[
f(y)\le f(x)+\inner{\nabla f(x)}{y-x}+\frac{L}{2}\norm{y-x}^2 .
\tag{6}
\]
\end{lemma}
\begin{lemma}[Three‑Point Identity]\label{lem:three}
For any $a,b\in\RR^d$ and any $c\in\RR^d$,
\[
\norm{a-c}^2=\norm{a-b}^2+2\inner{a-b}{b-c}+\norm{b-c}^2 .
\tag{7}
\]
\end{lemma}
\begin{lemma}[Noise Moment]\label{lem:noise}
Let $\xi\sim\mathcal N(0,I_d)$. Then $\EE[\xi]=0$ and $\EE[\norm{\xi}^2]=d$.
\end{lemma}

\section*{Main Proof (Step‑by‑Step Derivation)}
We prove Theorem \ref{thm:main}. All expectations are taken w.r.t. the randomness up to iteration $k$ unless otherwise noted.

\bigskip
\noindent\textbf{[Step 1] Distance Recursion.}  
From the update \eqref{1} and Lemma \ref{lem:three} with $a=x_{k+1}$, $b=x_k$, $c=x^\star$,
\[
\begin{aligned}
\norm{x_{k+1}-x^\star}^2
&=\norm{x_k-x^\star}^2
-2\alpha\inner{\nabla f(x_k)}{x_k-x^\star}
+2\sqrt{2\alpha\tau_k}\inner{\xi_k}{x_k-x^\star}   \\
&\quad+\alpha^2\norm{\nabla f(x_k)}^2
+2\alpha\sqrt{2\alpha\tau_k}\inner{-\nabla f(x_k)}{\xi_k}   \\
&\quad+2\alpha\tau_k\norm{\xi_k}^2 .
\end{aligned}
\tag{8}
\]

\noindent\textbf{[Step 2] Conditional Expectation.}  
Take $\EE[\cdot\mid\mathcal F_k]$; using $\EE[\xi_k\mid\mathcal F_k]=0$ and $\EE[\inner{-\nabla f(x_k)}{\xi_k}\mid\mathcal F_k]=0$,
\[
\EE\!\bigl[\norm{x_{k+1}-x^\star}^2\mid\mathcal F_k\bigr]
=
\norm{x_k-x^\star}^2
-2\alpha\inner{\nabla f(x_k)}{x_k-x^\star}
+\alpha^2\norm{\nabla f(x_k)}^2
+2\alpha\tau_k\,\EE\!\bigl[\norm{\xi_k}^2\bigr].
\tag{9}
\]

\noindent\textbf{[Step 3] Noise Moment.}  
By Lemma \ref{lem:noise}, $\EE[\norm{\xi_k}^2]=d$. Hence
\[
\EE\!\bigl[\norm{x_{k+1}-x^\star}^2\mid\mathcal F_k\bigr]
=
\norm{x_k-x^\star}^2
-2\alpha\inner{\nabla f(x_k)}{x_k-x^\star}
+\alpha^2\norm{\nabla f(x_k)}^2
+2\alpha d\,\tau_k .
\tag{10}
\]

\noindent\textbf{[Step 4] Use Convexity.}  
Convexity of $f$ yields
\[
f(x_k)-f^\star\le\inner{\nabla f(x_k)}{x_k-x^\star}.
\tag{11}
\]

\noindent\textbf{[Step 5] Use Smoothness to bound $\norm{\nabla f(x_k)}^2$.}  
From $L$‑smoothness (Lemma \ref{lem:smooth}) applied with $y=x^\star$,
\[
f^\star\ge f(x_k)+\inner{\nabla f(x_k)}{x^\star-x_k}
-\frac{L}{2}\norm{x_k-x^\star}^2 .
\]
Rearranging gives
\[
\inner{\nabla f(x_k)}{x_k-x^\star}\ge f(x_k)-f^\star-\frac{L}{2}\norm{x_k-x^\star}^2 .
\tag{12}
\]
Moreover, the co‑coercivity of gradients for $L$‑smooth convex functions implies
\[
\norm{\nabla f(x_k)}^2\le 2L\bigl(f(x_k)-f^\star\bigr).
\tag{13}
\]

\noindent\textbf{[Step 6] Plug (12)–(13) into (10).}  
\[
\begin{aligned}
\EE\!\bigl[\norm{x_{k+1}-x^\star}^2\mid\mathcal F_k\bigr]
&\le \norm{x_k-x^\star}^2
-2\alpha\Bigl(f(x_k)-f^\star-\frac{L}{2}\norm{x_k-x^\star}^2\Bigr)   \\
&\quad+ \alpha^2\cdot 2L\bigl(f(x_k)-f^\star\bigr)
+2\alpha d\,\tau_k .
\end{aligned}
\tag{14}
\]

\noindent\textbf{[Step 7] Collect terms in $f(x_k)-f^\star$.}  
Since $\alpha\le 1/L$, we have $1- \alpha L\ge 0$. Hence
\[
\begin{aligned}
\EE\!\bigl[\norm{x_{k+1}-x^\star}^2\mid\mathcal F_k\bigr]
&\le \norm{x_k-x^\star}^2
-2\alpha\bigl(1-\alpha L\bigr)\bigl(f(x_k)-f^\star\bigr)   \\
&\quad+ \alpha L \norm{x_k-x^\star}^2
+2\alpha d\,\tau_k .
\end{aligned}
\tag{15}
\]

\noindent\textbf{[Step 8] Rearranging.}  
Bring the term $\alpha L \norm{x_k-x^\star}^2$ to the left side:
\[
\EE\!\bigl[\norm{x_{k+1}-x^\star}^2\mid\mathcal F_k\bigr]
\le (1-\alpha L)\norm{x_k-x^\star}^2
-2\alpha\bigl(1-\alpha L\bigr)\bigl(f(x_k)-f^\star\bigr)
+2\alpha d\,\tau_k .
\tag{16}
\]

\noindent\textbf{[Step 9] Take full expectation.}  
Denote $\Delta_k:=\EE\!\bigl[f(x_k)-f^\star\bigr]$ and $R_k:=\EE\!\bigl[\norm{x_k-x^\star}^2\bigr]$.  
Taking expectation of \eqref{16} gives
\[
R_{k+1}\le (1-\alpha L)R_k-2\alpha(1-\alpha L)\Delta_k+2\alpha d\,\tau_k .
\tag{17}
\]

\noindent\textbf{[Step 10] Telescoping sum.}  
Rearrange \eqref{17} to isolate $\Delta_k$:
\[
\Delta_k\le \frac{R_k-R_{k+1}}{2\alpha(1-\alpha L)}+\frac{d\,\tau_k}{1-\alpha L}.
\tag{18}
\]
Summing from $k=0$ to $K-1$,
\[
\sum_{k=0}^{K-1}\Delta_k
\le \frac{R_0-R_{K}}{2\alpha(1-\alpha L)}+\frac{d}{1-\alpha L}\sum_{k=0}^{K-1}\tau_k .
\tag{19}
\]

\noindent\textbf{[Step 11] Upper bound $R_K$.}  
Since $R_K\ge 0$, drop it:
\[
\sum_{k=0}^{K-1}\Delta_k
\le \frac{R_0}{2\alpha(1-\alpha L)}+\frac{d}{1-\alpha L}\sum_{k=0}^{K-1}\tau_k .
\tag{20}
\]

\noindent\textbf{[Step 12] Use the specific schedule \eqref{2}.}  
Because $\tau_k=\tau_0/(k+1)$,
\[
\sum_{k=0}^{K-1}\tau_k
=\tau_0\sum_{k=1}^{K}\frac{1}{k}
\le \tau_0\bigl(1+\log K\bigr)
\le \tau_0\log (K+1)+\tau_0 .
\tag{21}
\]

\noindent\textbf{[Step 13] Relate average suboptimality to the last iterate.}  
Convexity of $f$ implies $f\bigl(\frac{1}{K}\sum_{k=0}^{K-1}x_k\bigr)\le\frac{1}{K}\sum_{k=0}^{K-1}f(x_k)$. Hence
\[
\EE\!\bigl[f(\bar x_K)\bigr]-f^\star
\le \frac{1}{K}\sum_{k=0}^{K-1}\Delta_k,
\qquad 
\bar x_K:=\frac{1}{K}\sum_{k=0}^{K-1}x_k .
\tag{22}
\]
Dividing \eqref{20} by $K$ and using \eqref{21} yields
\[
\EE\!\bigl[f(\bar x_K)\bigr]-f^\star
\le
\frac{R_0}{2\alpha(1-\alpha L)K}
+\frac{d\,\tau_0}{(1-\alpha L)}\frac{\log (K+1)}{K}
+\frac{d\,\tau_0}{(1-\alpha L)K}.
\tag{23}
\]

\noindent\textbf{[Step 14] Simplify constants.}  
Recall $R_0=\norm{x_0-x^\star}^2$ and $1-\alpha L\ge \tfrac12$ for $\alpha\le \tfrac{1}{2L}$ (we may tighten the bound by choosing $\alpha\le 1/L$). Replacing $1-\alpha L$ by a constant $c\in(0,1]$ gives the compact form
\[
\EE\!\bigl[f(\bar x_K)\bigr]-f^\star
\le
\frac{\norm{x_0-x^\star}^2}{2\alpha K}
+\alpha L\,\tau_0\frac{\log (K+1)}{K}
+ \frac{\alpha L\,\tau_0}{K}.
\tag{24}
\]
Dropping the last $O(1/K)$ term leads to the statement of Theorem \ref{thm:main} (inequality \eqref{5}). $\square$

\section*{Rate Derivation (With Constants)}
From \eqref{5},
\[
\EE\!\bigl[f(\bar x_K)-f^\star\bigr]
\le
\underbrace{\frac{\norm{x_0-x^\star}^2}{2\alpha}}_{\text{deterministic GD term}}\frac{1}{K}
\;+\;
\underbrace{\alpha L\tau_0}_{\text{noise‑bias term}}\frac{\log (K+1)}{K}.
\]
Choosing $\alpha = \Theta\!\bigl(1/\sqrt{K}\bigr)$ balances the two contributions and yields
\[
\EE\!\bigl[f(\bar x_K)-f^\star\bigr]=O\!\Bigl(\frac{\log K}{\sqrt{K}}\Bigr),
\]
whereas any fixed $\alpha\le 1/L$ gives the $O\bigl((\log K)/K\bigr)$ rate stated in Theorem \ref{thm:main}.

\section*{Special Cases}
\begin{enumerate}[label=\alph*.]
\item \textbf{Zero Noise ($\tau_0=0$).}  The bound collapses to $\frac{\norm{x_0-x^\star}^2}{2\alpha K}$, i.e. the classical GD rate.
\item \textbf{Constant Temperature ($\tau_k\equiv\tau$).}  The sum $\sum_{k=0}^{K-1}\tau_k = K\tau$ and \eqref{23} becomes 
\[
\EE[f(\bar x_K)]-f^\star\le \frac{\norm{x_0-x^\star}^2}{2\alpha K}+2\alpha d\,\tau,
\]
exhibiting a non‑vanishing bias term $2\alpha d\,\tau$.
\item \textbf{Faster Decay $\tau_k=\tau_0/(k+1)^{1+\epsilon}$, $\epsilon>0$.}  
Then $\sum_{k=0}^{K-1}\tau_k = O(1)$, and the logarithmic factor disappears, yielding the pure $O(1/K)$ rate.
\end{enumerate}

\end{document}